\bookchapter{الفصل 58: دليل المراجعة السريعة والمصادر}

\chapterintro{هذا الفصل دليل مراجعة سريعة صالح للاستخدام المستقل. عندما تخلط عليك طريقة كتابة ما أو تحتاج إلى قرار سريع، افتحه مباشرة. ننصحك بجعله «بطاقة جيب»: مراجعة الكود الشائع، والأسئلة المتكررة، ومصادر التعلم، دون تفويت صفحة.}

\booksection{أولاً: بطاقة المراجعة السريعة للكود الشائع}

\begin{center}
\begin{tabularx}{\linewidth}{|X|X|}
\hline
\textbf{ماذا تريد أن تفعل} & \textbf{الكود الأساسي} \tabularnewline
\hline
حفظ حالة قابلة للتغيير & var x by remember \{ mutableStateOf(القيمة الابتدائية) \} \tabularnewline
عدم فقدان الحالة عند تدوير الشاشة & rememberSaveable \{ mutableStateOf(القيمة الابتدائية) \} \tabularnewline
تخطيط عمودي/أفقي & \texttt{Column \{ \}}/\texttt{Row \{ \}} \tabularnewline
قائمة بالتحميل الكسول & LazyColumn \{ items(القائمة) \{ ... \} \} \tabularnewline
التنقل بين الشاشات & navController.navigate("المسار") \tabularnewline
تنفيذ أعمال الانتظار داخل coroutine & \texttt{viewModelScope.\allowbreak launch \{ .\allowbreak .\allowbreak .\allowbreak  \}} \tabularnewline
مراقبة تدفق البيانات & \texttt{val list by flow.\allowbreak collectAsState()} \tabularnewline
تخزين إعدادات صغيرة & context.dataStore.edit \{ it[KEY] = القيمة \} \tabularnewline
تخزين بيانات مهيكلة & ثلاثية Room: @Entity / @Dao / @Database \tabularnewline
جلب البيانات عبر الإنترنت & واجهة Retrofit + دالة suspend \tabularnewline
إرسال إشعار & NotificationChannel + NotificationCompat.Builder \tabularnewline
كتابة السجلات (Log) & \texttt{Log.\allowbreak d("TAG", "message")} \tabularnewline
\hline
\end{tabularx}
\end{center}

\begin{center}\small\color{gray}بطاقة مراجعة تطوير Android (إصدار أغسطس 2026)\end{center}

\booksection{ثانياً: الأسئلة الشائعة}

س1: لا أعرف البرمجة إطلاقاً، فكم أحتاج من الوقت لتعلم تطوير Android؟

اتبع مسار هذا الكتاب: في الأسبوعين 1-2 جهّز البيئة وتعلم أساسيات Kotlin؛ في الأسبوعين 3-4 تعلم Compose وابنِ واجهة؛ في الأسبوعين 5-6 أكمل مشروع دفتر المصاريف؛ في الأسبوعين 7-8 أتم مشروع الطقس. شهران تنتقل فيهما من الصفر إلى القدرة على صنع تطبيق صغير باستقلال. الشرط أن تخصص 1-2 ساعة يومياً وأن تكتب الكود بيدك لا أن تكتفي بالقراءة.

س2: هل ما زلت بحاجة إلى تعلم Java؟

لا داعي. المشاريع الجديدة والدروس الجديدة والوظائف الجديدة تعتمد Kotlin أساساً. ولن تصطدم بكثير من Java إلا عند «صيانة مشروع عمره عشر سنوات». وبعد إتقان Kotlin ستعوّض Java سريعاً إن اقتضت العمل — فكلاهما يقوم على بيئة الجهاز الافتراضي نفسها.

س3: المحاكي أم الجهاز الحقيقي؟

المحاكي مريح للتطوير والتصحيح (لقطات الحالة وتعدد الطرازات)، لكن القبول النهائي لا بد أن يكون على جهاز حقيقي — فاللمس والأداء والشبكة الفعلية والمستشعرات لا يعوّضها المحاكي. وإن سمحت الظروف فاقتناء جهاز حقيقي متوسط الفئة للاختبار هو أجدى استثمار.

س4: هل سوق العمل جيد لتطوير Android؟

أوضاع 2026: المنافسة شديدة على الوظائف المبتدئة التي تقتصر على «كتابة الواجهات»، لكن المطوّر الذي «يتقن Kotlin + يفهم البنية + لديه عمل منشور» ما زال مطلوباً بشدة. وأنت في فترة التعلم، أفضل سيرة ذاتية لك هو ذلك التطبيق الذي صنعته بنفسك في المتجر.

س5: ماذا أفعل إذا علقت أمام رسالة خطأ؟

الخطوات الذهبية الثلاث: (1) ابحث عن السطر الأول من الخطأ حرفياً في StackOverflow أو محرك بحث؛ (2) راجع الصفحة المقابلة في الوثائق الرسمية؛ (3) إن لم ينجح ذلك، اجمع «السلوك المتوقع + الظاهر الفعلي + نص الخطأ كاملاً + الكود ذي الصلة» وأرسلها سؤالاً إلى الذكاء الاصطناعي أو المجتمع. الغالبية العظمى من المشاكل تُحل في الخطوة الأولى.

\booksection{ثالثاً: مصادر تعلم موصى بها}

\begin{center}
\begin{tabularx}{\linewidth}{|X|X|X|}
\hline
\textbf{النوع} & \textbf{المصدر} & \textbf{سبب التوصية} \tabularnewline
\hline
رسمي & developer.android.com & الوثائق الرسمية لـ Android، وفيها دروس Compose وKotlin — المرجع الأول والأوثق \tabularnewline
رسمي & موقع Kotlin الرسمي (kotlinlang.org) & وثائق اللغة الرسمية + ساحة تجريب على الإنترنت \tabularnewline
رسمي & Google Codelabs & دروس عملية خطوة بخطوة يد بيد، وأشد ودّية للمبتدئين \tabularnewline
مجتمع & Stack Overflow (وسم android) & المحطة الأولى لطلب المساعدة عند الأخطاء \tabularnewline
مجتمع & قسما Android في Juejin وCSDN & مقالات عملية بالصينية كثيرة؛ انتبه للتمييز بين الجديد والقديم (اقرأ التاريخ أولاً) \tabularnewline
فيديو & قناة Google Developers الرسمية & شرح التقنيات الجديدة في مؤتمر المطورين كل عام \tabularnewline
مرافق للكتاب & القسم السابع «التطبيق العملي لبرمجة الذكاء الاصطناعي» (الفصول 37-47) & المنهجية الكاملة لتطوير Android بمساعدة الذكاء الاصطناعي \tabularnewline
\hline
\end{tabularx}
\end{center}

\begin{center}\small\color{gray}قائمة مصادر تعلم تطوير Android (أغسطس 2026)\end{center}

\booksection{رابعاً: مراجعة سريعة لأوامر Gradle الشائعة}

النقر على الأزرار يكفي عادة، لكن سطر الأوامر أجدى أحياناً (خصوصاً مع أدوات الذكاء الاصطناعي أو خطوط CI):

\begin{center}
\begin{tabularx}{\linewidth}{|X|X|}
\hline
\textbf{الأمر} & \textbf{الوظيفة} \tabularnewline
\hline
\texttt{.\allowbreak /\allowbreak gradlew assembleDebug} & بناء APK نسخة التصحيح (على Windows احذف ./) \tabularnewline
\texttt{.\allowbreak /\allowbreak gradlew assembleRelease} & بناء نسخة النشر (يتطلب ضبط التوقيع) \tabularnewline
\texttt{.\allowbreak /\allowbreak gradlew test} & تشغيل جميع اختبارات الوحدة \tabularnewline
\texttt{.\allowbreak /\allowbreak gradlew lint} & فحص ساكن: كشف الأخطاء الكامنة والكتابات غير القياسية \tabularnewline
\texttt{.\allowbreak /\allowbreak gradlew clean} & تنظيف مخلفات البناء (النظير التصريفي لأسلوب «إعادة التشغيل») \tabularnewline
\texttt{.\allowbreak /\allowbreak gradlew :app: dependencies} & عرض شجرة التبعيات (سلاح ممتاز لتشخيص تعارض الإصدارات) \tabularnewline
\hline
\end{tabularx}
\end{center}

\begin{center}\small\color{gray}مراجعة سريعة لأوامر Gradle الشائعة\end{center}

\begin{tcbinfo}{نصيحة: عند مشكلة في البناء ابدأ بـ clean}
ظاهرة «كان يعمل أمس وتوقف اليوم» سببها في ست أو سبع حالات من كل عشر خلل في ذاكرة البناء المؤقتة. بعد ./gradlew clean أعد البناء، فذلك يشفي أغلب المشاكل الغامضة. وإن لم ينجح فأعد تشغيل Android Studio — وخذ هذه الحيلة مع File Invalidate Caches معاً.

\end{tcbinfo}

\booksection{خامساً: خاتمة — من أول تطبيق إلى التطبيق رقم N}

\booksection{سابعاً: خاتمة — من أول تطبيق إلى التطبيق رقم N}

بالنظر إلى الوراء، انتقلت من «لا أعرف ما هو Activity» إلى تجهيز البيئة وإتقان Kotlin ورسم الواجهات بـ Compose وربط التطبيق بقاعدة بيانات وشبكة، ثم الاختبار والتجزئة والنشر — في هذا الطريق قطعتَ أصعب جزء من الرحلة من الصفر إلى الواحد.

عالم التقنية يتغير بسرعة: الأطر تترقى، والأدوات تتبدل، والذكاء الاصطناعي يتولى المزيد من أعمال «ضرب المفاتيح». لكن شيئاً واحداً لا يبطل: قدرتك على تفكيك فكرة غامضة إلى خطوات واضحة ثم تنفيذها خطوة خطوة. الأدوات مضخّمات، والفكر أساس. المضخّم يُستبدل، والأساس لا يهتز.

لم تعد مجرد "شخص يعرف استخدام الهاتف". يمكنك أن تحوّل أفكارك في رأسك إلى تطبيقات تعمل فعلًا على ثلاثمئة مليون جهاز أندرويد. من اليوم، ابتكر بلوحة المفاتيح — هذه هي رومانسية مطوّري أندرويد.

\booksection{المراجع}

\begin{enumerate}
  \item Android Developers. (2026). Android Developer Guides and Jetpack Compose Documentation. developer.android.com.
  \item Kotlin Foundation. (2026). Kotlin Language Documentation. kotlinlang.org.
  \item Android Open Source Project. (2026). Android Platform Version History and API Levels.
  \item Google Play. (2026). Play Console Help: Publishing and App Signing.
  \item AndroidX Release Notes. (2026). Room / DataStore / Navigation / WorkManager releases.
  \item Square. (2026). Retrofit Documentation. github.com/square/retrofit.
\end{enumerate}
